\section{Sonuç}
\label{sec:conclusion}

Bu çalışmada, HiperKvasir 23-sınıf GI endoskopi sınıflandırması için çoklu ViT
tabanlı bir özellik füzyonu yaklaşımı sunulmuştur. Üç tamamlayıcı omurganın
(ViT-B/16, Swin-T, BEiT-B/16) ağırlıklı füzyonunun ince ayarlı sürümü, resmi
5-katlı protokolde en iyi yapılandırma olmuş; havuzlanmış makro-F1
$0{,}6119$ (\%95 GA $[0{,}5930,\,0{,}6295]$), doğruluk $\approx 0{,}88$
elde edilmiştir. Bu yapılandırma, eşli McNemar testinde hem en iyi tekli
omurgadan hem de en iyi ikili füzyondan istatistiksel olarak anlamlı biçimde
daha doğrudur.

Çalışmanın temel çıkarımları şunlardır: (i)~çok-omurga füzyonu genel doğruluğa
anlamlı bir katkı sağlar, ancak makro-F1 kazancı nadir sınıflarca sınırlanır;
(ii)~ağırlıklı füzyon, birleştirme ve gelişmiş GMU füzyonundan daha iyi/eşdeğer
sonuç vermiştir; (iii)~son blokların ince ayarı dondurulmuş öznitelik çıkarımını
geçmiştir;
(iv)~dengesizlik için ağırlıklı örnekleyici yeterli olup son-işlem logit
düzeltmesi başarımı düşürmüştür. TTA ve seed topluluğu gibi toplamsal teknikler
güven aralığı düzeyinde ayırt edilebilir bir iyileşme getirmemiştir; bu, mimarinin
bu protokolde tavanına ulaştığını gösterir.

\paragraph{Sınırlamalar ve gelecek çalışmalar.} Temel sınırlama, çok düşük
destekli nadir sınıflardan kaynaklanan bir veri kıtlığıdır. Gelecekte;
(a)~ordinal yapıyı (ör.\ kolit dereceleri) açıkça modelleyen sıralı
sınıflandırma kayıpları, (b)~nadir sınıflar için hedefli veri toplama, üretici
artırma veya odak kaybı gibi dengesizlik-duyarlı kayıplar~\cite{lin2017focal}, ve
(c)~klinik açıdan ``yakın'' hataları ayrıştıracak ince taneli ayrım öğrenme
yaklaşımları umut verici yönlerdir. Bildirilen sayıların tümü
kaydedilen değerlendirme çıktılarına dayanır; literatür karşılaştırmaları
protokol farkları nedeniyle bağlamsaldır ve güncel-en-iyi bir başarı iddiasında
bulunulmamıştır.

\paragraph{Proje deposu.} Kaynak kod, eğitim çıktıları ve tanıtım videosu
bağlantısı proje deposunda yer alır:
\url{https://github.com/YasinEkici/hyperkvasir-multi-backbone-fusion}.
